%% ma: L12-B13-0045
%% lop: 12
%% chuong: 4
%% bai: 13
%% dang: 12.13.24
%% loai: TL
%% muc: VD
%% dap_an: "-"
%% nguon: "(NBV)-10 ĐỀ KIỂM TRA HK2-TOÁN 12-DUNG CHUNG 3 SÁCH.pdf (H04, Đề số 4), Phần 4 Câu 1 (tự luận)"
%% kien_thuc: "Bài 13 (thể tích khối tròn xoay giới hạn bởi hai đồ thị, quay quanh trục Ox)"
%% trang_thai: chinh-thuc
%% ghi_chu: "Đáp án gốc V = pi(19/3 - 3/ln4) khớp. Hình gốc tô miền dưới đồ thị đến trục $Ox$ (không phải miền (H) của đề); hình kho dựng lại, tô đúng miền (H) giữa y=2^x, y=3-x và trục Oy. Đã duyệt 24/10/2026."
\begin{ex}
Gọi $(H)$ là hình phẳng giới hạn bởi $(C):y=2^x$, $(d):y=-x+a$ và trục $Oy$. Biết rằng $(C)$ và $(d)$ cắt nhau tại một điểm duy nhất có hoành độ bằng $1$. Tính thể tích $V$ của khối tròn xoay tạo bởi $(H)$ khi nó quay quanh trục $Ox$.
\begin{center}
\begin{tikzpicture}[xscale=1.5,yscale=.9]
  \fill[gray!45] plot[domain=0:1,samples=40] (\x,{2^\x}) -- (1,2) -- (0,3) -- cycle;
  \draw[thick,domain=-0.9:1.45,samples=60] plot (\x,{2^\x});
  \draw[thick] (-0.3,3.3)--(3.2,-0.2);
  \draw[truc] (-1,0)--(3.5,0) node[below]{$x$};
  \draw[truc] (0,-0.4)--(0,3.8) node[left]{$y$};
  \draw[phu] (1,0)--(1,2);
  \path (0,0) node[below left]{\footnotesize$O$} (1,0) node[below]{\footnotesize$1$} (3,0) node[below]{\footnotesize$3$};
  \path (2.4,0.75) node[above right]{\footnotesize$(d)$} (1.4,2.64) node[right]{\footnotesize$(C)$} (0.42,2.05) node{\footnotesize$(H)$};
\end{tikzpicture}
\end{center}
\loigiai{Hoành độ giao điểm bằng $1$ nên $2^1=-1+a$, suy ra $a=3$ và $(d):y=3-x$.

Trên $[0;1]$ có $3-x\ge2^x>0$, nên khối tròn xoay có thể tích
\[V=\pi\int_0^1\left[(3-x)^2-4^x\right]\mathrm dx=\pi\left[-\dfrac{(3-x)^3}3-\dfrac{4^x}{\ln4}\right]_0^1=\pi\left(\dfrac{19}3-\dfrac3{\ln4}\right).\]}
\end{ex}
